\begin{abstract}
We train a small transformer on a non-ergodic dataset consisting of Mess3 processes with different parameters, where each sequence is generated by a single ergodic component. The transformer must solve a hierarchical inference problem: identifying the generating component (meta-synchronization) and tracking belief states within it. We derive mathematical predictions for the expected activation geometry, connecting the non-ergodic structure to the Factored World Hypothesis. We first verify that a single Mess3 process produces Sierpinski gasket geometry in the residual stream ($R^2 = 0.989$), then show that the non-ergodic mixture induces a hierarchical representation where component identity and within-component beliefs are encoded in approximately orthogonal subspaces. We analyze synchronization dynamics---how the model identifies the active component across context positions and layers---revealing that meta-synchronization precedes belief refinement. This setup provides a minimal, analytically tractable model of in-context learning as source discrimination.
\end{abstract}
